\documentclass[12pt,a4paper]{article}

\usepackage[margin=2.2cm]{geometry}
\usepackage{xcolor}
\usepackage{mdframed}
\usepackage{listings}
\usepackage{titlesec}
\usepackage{enumitem}
\usepackage{booktabs}
\usepackage{array}
\usepackage{graphicx}
\usepackage{fancyhdr}
\usepackage{tikz}
\usepackage{parskip}
\usepackage[hidelinks]{hyperref}
\usepackage{fontenc}

% ── Colors ────────────────────────────────────────────────────────
\definecolor{navyblue}{HTML}{083C92}
\definecolor{lightblue}{HTML}{EBF2FA}
\definecolor{safegreen}{HTML}{1A6B3C}
\definecolor{lightgreen}{HTML}{E8F5EE}
\definecolor{unsafered}{HTML}{8B1A1A}
\definecolor{lightred}{HTML}{FDECEA}
\definecolor{reducegold}{HTML}{B8860B}
\definecolor{lightgold}{HTML}{FDF6E3}
\definecolor{codegray}{HTML}{F4F4F4}
\definecolor{darkgray}{HTML}{444444}

% ── Section styling ───────────────────────────────────────────────
\titleformat{\section}
  {\Large\bfseries\color{navyblue}}
  {\thesection.}{0.5em}{}
  [\color{navyblue}\rule{\linewidth}{1.5pt}]

\titleformat{\subsection}
  {\large\bfseries\color{darkgray}}
  {\thesubsection.}{0.5em}{}

% ── Code style ────────────────────────────────────────────────────
\lstset{
  basicstyle=\ttfamily\small,
  backgroundcolor=\color{codegray},
  frame=single,
  framerule=0pt,
  rulecolor=\color{codegray},
  breaklines=true,
  captionpos=b,
  numbers=none,
  xleftmargin=10pt,
  xrightmargin=10pt,
}

% ── Coloured boxes ────────────────────────────────────────────────
\newmdenv[
  backgroundcolor=lightgreen,
  linecolor=safegreen,
  linewidth=2pt,
  leftline=true, rightline=false, topline=false, bottomline=false,
  innerleftmargin=10pt, innerrightmargin=10pt,
  innertopmargin=8pt, innerbottommargin=8pt,
  skipabove=6pt, skipbelow=6pt,
]{safebox}

\newmdenv[
  backgroundcolor=lightred,
  linecolor=unsafered,
  linewidth=2pt,
  leftline=true, rightline=false, topline=false, bottomline=false,
  innerleftmargin=10pt, innerrightmargin=10pt,
  innertopmargin=8pt, innerbottommargin=8pt,
  skipabove=6pt, skipbelow=6pt,
]{unsafebox}

\newmdenv[
  backgroundcolor=lightgold,
  linecolor=reducegold,
  linewidth=2pt,
  leftline=true, rightline=false, topline=false, bottomline=false,
  innerleftmargin=10pt, innerrightmargin=10pt,
  innertopmargin=8pt, innerbottommargin=8pt,
  skipabove=6pt, skipbelow=6pt,
]{reducebox}

\newmdenv[
  backgroundcolor=lightblue,
  linecolor=navyblue,
  linewidth=2pt,
  leftline=true, rightline=false, topline=false, bottomline=false,
  innerleftmargin=10pt, innerrightmargin=10pt,
  innertopmargin=8pt, innerbottommargin=8pt,
  skipabove=6pt, skipbelow=6pt,
]{infobox}

% ── Header / Footer ───────────────────────────────────────────────
\pagestyle{fancy}
\fancyhf{}
\fancyhead[L]{\small\color{navyblue}\textbf{FlangParallelAnalyzer}}
\fancyhead[R]{\small\color{darkgray}RV College of Engineering}
\fancyfoot[C]{\small\thepage}
\renewcommand{\headrulewidth}{0.5pt}
\renewcommand{\headrule}{\color{navyblue}\hrule width\headwidth height\headrulewidth}

% =================================================================
\begin{document}

% ── Title page ───────────────────────────────────────────────────
\begin{titlepage}
\centering
\vspace*{1.5cm}
{\color{navyblue}\rule{\linewidth}{3pt}}\\[1.2cm]
{\Huge\bfseries\color{navyblue} FlangParallelAnalyzer}\\[0.6cm]
{\Large\color{darkgray} A Simple Guide to the Project}\\[0.4cm]
{\large\color{darkgray}\textit{Automatic Loop Parallelism Analysis for Fortran}}\\[1.5cm]
{\color{navyblue}\rule{\linewidth}{1pt}}\\[1.5cm]

\begin{tabular}{c}
{\large Aryan Gupta} \\[4pt]
{\large Arushi Vaidya} \\[4pt]
{\large Arpita} \\[12pt]
{\normalsize Department of Computer Science and Engineering} \\[4pt]
{\normalsize RV College of Engineering, Bengaluru} \\
\end{tabular}

\vfill
{\color{navyblue}\rule{\linewidth}{3pt}}\\[0.5cm]
{\normalsize 2024--25}
\end{titlepage}

\tableofcontents
\newpage

% =================================================================
\section{What is this project about?}

Imagine you have a Fortran program that processes a huge array of
numbers --- say, a weather simulation running over a million grid
points. Inside the program, there is a loop that looks like this:

\begin{lstlisting}[language=Fortran]
do i = 1, 1000000
    temperature(i) = temperature(i) * 1.05
end do
\end{lstlisting}

Your computer today has 8, 16, or even 64 cores sitting idle. Instead
of running this loop one step at a time, you could split it across all
cores and finish it 64$\times$ faster. That is what \textbf{parallel
computing} is --- doing many things at the same time.

In Fortran, you tell the compiler to parallelize a loop by adding a
single line:

\begin{lstlisting}[language=Fortran]
!$OMP PARALLEL DO          <- this one line splits the work
do i = 1, 1000000
    temperature(i) = temperature(i) * 1.05
end do
\end{lstlisting}

\begin{infobox}
\textbf{The problem:} Not every loop can be parallelized safely.
If one iteration of the loop depends on a result computed by a previous
iteration, running them at the same time gives the \textbf{wrong
answer} --- silently, with no error message.

Figuring out which loops are safe to parallelize, by hand, for a
codebase with hundreds of loops, takes hours and is error-prone.
\end{infobox}

\textbf{FlangParallelAnalyzer (FPA)} solves this automatically. You
give it a Fortran file, and it tells you --- for every loop ---
whether it is safe to parallelize, needs a special clause, or is
genuinely unsafe.

% =================================================================
\section{The Three Possible Outcomes}

For every loop it finds, FPA gives one of three verdicts:

\vspace{6pt}
\begin{safebox}
\textbf{\color{safegreen}SAFE --- Go ahead and parallelize.}\\
Each iteration of the loop works on completely independent data.
No two iterations ever read or write the same memory cell.\\[4pt]
\textbf{Example:}
\begin{lstlisting}[language=Fortran]
do i = 1, n
    b(i) = a(i) * 2.0     ! iteration 3 touches b(3) and a(3)
end do                     ! iteration 7 touches b(7) and a(7)
                           ! they never overlap -> SAFE
\end{lstlisting}
\textbf{Hint emitted:} \texttt{!\$OMP PARALLEL DO}
\end{safebox}

\vspace{6pt}
\begin{reducebox}
\textbf{\color{reducegold}REDUCTION --- Safe with a special clause.}\\
The loop adds up (or multiplies) into a single accumulator variable.
OpenMP can handle this pattern natively.\\[4pt]
\textbf{Example:}
\begin{lstlisting}[language=Fortran]
do i = 1, n
    total = total + a(i) * b(i)   ! "total" accumulates across iters
end do
\end{lstlisting}
\textbf{Hint emitted:} \texttt{!\$OMP PARALLEL DO REDUCTION(+:total)}
\end{reducebox}

\vspace{6pt}
\begin{unsafebox}
\textbf{\color{unsafered}UNSAFE --- Do not parallelize.}\\
One iteration reads a value that was written by another iteration.
Parallelizing this would produce wrong results.\\[4pt]
\textbf{Example:}
\begin{lstlisting}[language=Fortran]
do i = 2, n
    a(i) = a(i) + a(i-1)   ! iteration 5 reads a(4), which
end do                      ! iteration 4 is still computing -> UNSAFE
\end{lstlisting}
\textbf{Hint emitted:} \texttt{! Cannot parallelize}
\end{unsafebox}

% =================================================================
\section{How Does It Work? (The Pipeline)}

FPA does not guess. It reads the compiled form of the Fortran program
(called \textbf{FIR} --- Fortran Intermediate Representation) and runs
a five-step analysis on every loop it finds.

\begin{infobox}
\textbf{Why FIR and not the Fortran source directly?}\\
Fortran source code has many tricky syntax rules (implicit variable
types, fixed vs free format, etc.). FIR is a clean, structured version
produced by the compiler where every array access, every variable, and
every loop is spelled out explicitly. It is much easier to analyze.
\end{infobox}

\subsection{Phase 1 --- What does this loop look like?}

FPA first collects basic facts about the loop:
\begin{itemize}[leftmargin=1.5em]
  \item What are the loop bounds? (e.g., \texttt{1} to \texttt{n})
  \item How deep is the nesting? (is this inside another loop?)
  \item How many operations are in the body?
\end{itemize}
At this point, the verdict is still \textit{unknown}.

\subsection{Phase 2 --- What memory does this loop touch?}

FPA walks every read (\texttt{load}) and write (\texttt{store}) inside
the loop and builds a list of every variable it sees:

\begin{center}
\begin{tabular}{llll}
\toprule
\textbf{Variable} & \textbf{Read?} & \textbf{Written?} & \textbf{External?} \\
\midrule
\texttt{a} (array)   & Yes & No  & Yes (passed in) \\
\texttt{b} (array)   & No  & Yes & Yes (passed in) \\
\texttt{i} (counter) & Yes & Yes & No  (local)     \\
\bottomrule
\end{tabular}
\end{center}

\textit{External} means the variable came from outside the loop ---
i.e., it is a function argument or a variable declared before the loop.
These are the ones that matter for dependency analysis.

\subsection{Phase 3 --- Are array subscripts independent?}

This is the core of the analysis. FPA looks at how each array is
indexed and asks: \textbf{``Does iteration \texttt{i} touch the same
element as iteration \texttt{j}?''}

It traces every subscript back through the compiler IR to find out
whether it is:

\begin{itemize}[leftmargin=1.5em]
  \item \textbf{Exactly the loop variable \texttt{i}:} each iteration
        touches a unique element --- independent, no conflict.

  \begin{lstlisting}[language=Fortran]
  a(i)      <- subscript = i  -> safe, each iter unique
  \end{lstlisting}

  \item \textbf{The loop variable plus or minus a constant (\texttt{i$\pm$k}):}
        iteration \texttt{i} reads the element that iteration
        \texttt{i-k} will write --- that is a
        \textbf{loop-carried dependency} and the loop is immediately
        marked \textsc{Unsafe}.

  \begin{lstlisting}[language=Fortran]
  a(i-1)    <- subscript = i-1  -> iter i reads what iter i-1 writes
  \end{lstlisting}

  \item \textbf{Something else} (e.g., an indirect index like
        \texttt{a(idx(i))}): FPA cannot prove it is safe, so it passes
        to Phase 4.
\end{itemize}

\subsection{Phase 4 --- Is this a reduction?}

If Phase 3 left the verdict unknown (usually because a scalar variable
is both read and written), Phase 4 checks for the \textbf{reduction
pattern}:

\begin{lstlisting}[language=Fortran]
! The pattern FPA looks for:
old_value  = load(accumulator)          ! step 1: read
new_value  = old_value + something      ! step 2: arithmetic
             store(new_value, accumulator)  ! step 3: write back
\end{lstlisting}

If this exact chain is found, the loop is marked \textsc{Reduction}
and the correct \texttt{REDUCTION(+:var)} or \texttt{REDUCTION(*:var)}
clause is emitted.

\subsection{Phase 5 --- Final conservative decision}

Any loop still undecided after Phase 4 gets a fallback verdict:

\begin{itemize}[leftmargin=1.5em]
  \item If no external variable is written at all: the loop is
        \textsc{Safe} (a read-only loop cannot have dependencies).
  \item Otherwise: \textsc{Unsafe} --- FPA does not know enough to
        guarantee correctness, so it plays it safe.
\end{itemize}

\begin{infobox}
\textbf{Key design principle:} FPA would rather call a safe loop
``unsafe'' (you miss a speedup) than call an unsafe loop ``safe''
(you get wrong answers). \textbf{Zero false positives} is guaranteed.
\end{infobox}

% =================================================================
\section{A Full Worked Example}

Here is what FPA produces when run on a real Fortran file:

\textbf{Input:}
\begin{lstlisting}[language=Fortran]
subroutine scale(a, b, n)
    real, intent(in)  :: a(n)
    real, intent(out) :: b(n)
    integer :: i
    do i = 1, n
        b(i) = a(i) * 2.0
    end do
end subroutine
\end{lstlisting}

\textbf{Command:}
\begin{lstlisting}[language=bash]
flang-new -fc1 -emit-fir scale.f90 -o /tmp/scale.fir
./build/tools/fpa-tool/fpa-tool /tmp/scale.fir
\end{lstlisting}

\textbf{Output:}
\begin{lstlisting}
[FlangParallelAnalyzer] Function: _QPscale
------------------------------------------------------------
  Loop #1 @ scale.fir:6
  Bounds : [1 .. ? step 1]
  Access : ext-reads=1  ext-writes=1  ext-readwrites=0
           [R]  array  -- a
           [W]  array  -- b
  Status : SAFE
  Hint   : !$OMP PARALLEL DO
  Reason : Independent per-element access: each iteration reads
           a(i) and writes b(i) with no overlap across iterations.
------------------------------------------------------------
\end{lstlisting}

The developer can now confidently add \texttt{!\$OMP PARALLEL DO}
above the loop and get correct parallel execution.

% =================================================================
\section{What the Tool is Made Of}

\begin{center}
\begin{tabular}{p{5cm} p{2cm} p{7cm}}
\toprule
\textbf{File} & \textbf{Size} & \textbf{What it does} \\
\midrule
\texttt{LoopParallelAnalysis.cpp} & 471 lines &
  The main analysis engine. Runs all 5 phases on every loop. \\[4pt]
\texttt{AccessClassifier.cpp} & 165 lines &
  Walks through every memory read/write in a loop and records it. \\[4pt]
\texttt{main.cpp} (fpa-tool) & 71 lines &
  The command-line tool. Loads the file and runs the analysis. \\[4pt]
\texttt{report.py} & 1259 lines &
  Python script that runs the tool on all test files and generates the interactive HTML report. \\[4pt]
\texttt{run\_tests.py} & 250 lines &
  Automated test runner --- checks that all 35 test cases pass. \\
\bottomrule
\end{tabular}
\end{center}

\vspace{6pt}
\textbf{Technologies used:}
\begin{itemize}[leftmargin=1.5em, itemsep=2pt]
  \item \textbf{LLVM 18} --- the industry-standard compiler toolkit
        that FPA is built on top of.
  \item \textbf{MLIR} --- the part of LLVM that lets us write our own
        analysis pass and walk through code systematically.
  \item \textbf{Flang} --- the LLVM Fortran compiler that converts
        \texttt{.f90} source to FIR (the format FPA reads).
  \item \textbf{C++ 17} --- the language the analysis logic is written
        in.
  \item \textbf{Python 3} --- used for the test runner and HTML report
        generator.
\end{itemize}

% =================================================================
\section{Test Results}

FPA was tested on a suite of 35 handcrafted Fortran programs covering
every major loop pattern. Every single test gets the correct verdict.

\begin{center}
\begin{tabular}{lcc}
\toprule
\textbf{Category} & \textbf{Test cases} & \textbf{Correct} \\
\midrule
Embarrassingly parallel (SAFE)   & 13 & 13 \\
Scalar reductions (REDUCTION)    &  5 &  5 \\
Loop-carried dependencies (UNSAFE) & 10 & 10 \\
Complex / indirect patterns (UNSAFE conservatively) &  7 &  7 \\
\midrule
\textbf{Total} & \textbf{35} & \textbf{35 (100\%)} \\
\bottomrule
\end{tabular}
\end{center}

\begin{safebox}
\textbf{Zero false positives.} FPA never tells you a loop is safe when
it is not. The worst it can do is miss a parallelization opportunity ---
it never gives wrong answers.
\end{safebox}

% =================================================================
\section{The HTML Report}

Running \texttt{python3 scripts/report.py} generates a browser-based
report with three panels:

\begin{enumerate}[leftmargin=1.5em]
  \item \textbf{File Browser} --- lists all Fortran files with a
        colour-coded badge (\textcolor{safegreen}{SAFE} /
        \textcolor{reducegold}{REDUCTION} /
        \textcolor{unsafered}{UNSAFE}) for each loop.
  \item \textbf{Source View} --- shows the Fortran source with an
        inline chip on the \texttt{do} line indicating the verdict
        and a confidence score (e.g.\ HIGH 90\%).
  \item \textbf{Analysis Detail} --- shows exactly which of the 5
        phases fired, what memory accesses were found, and the full
        reason for the verdict.
\end{enumerate}

This lets a developer immediately see what was decided and \textit{why},
without reading any compiler IR.

% =================================================================
\section{Known Limitations}

FPA is a research prototype. A few patterns are not yet handled:

\begin{itemize}[leftmargin=1.5em, itemsep=4pt]
  \item \textbf{In-place same-index updates:} \texttt{a(i) = a(i) * 3}
        is actually safe (each iteration touches its own element) but
        FPA conservatively marks it \textsc{Unsafe} because it cannot
        confirm the read happens before the write without extra
        analysis.

  \item \textbf{Multi-dimensional arrays:} \texttt{c(i,j)} subscripts
        are not yet traced, so nested loops over 2D arrays get
        \textsc{Unsafe}.

  \item \textbf{Function calls:} if the loop body calls an external
        function, FPA does not know if that function has side effects,
        so it cannot guarantee safety.

  \item \textbf{Min/Max reductions:} only \texttt{+} and \texttt{*}
        reductions are detected. \texttt{MIN}/\texttt{MAX}
        accumulations are not yet supported.
\end{itemize}

% =================================================================
\section{What's Next}

\begin{itemize}[leftmargin=1.5em, itemsep=4pt]
  \item Fix the in-place update false negative with a read-before-write
        check on the value chain.
  \item Add support for \texttt{MIN}/\texttt{MAX} reductions.
  \item Trace multi-dimensional subscripts axis by axis.
  \item Analyze function side-effects interprocedurally.
  \item Automatically insert \texttt{!\$OMP} directives into the
        Fortran source (currently the tool only suggests them).
\end{itemize}

% =================================================================
\section{Summary}

\begin{infobox}
\textbf{In one paragraph:}

FlangParallelAnalyzer is a static analysis tool that reads compiled
Fortran programs and automatically determines, for every \texttt{DO}
loop, whether it is safe to run in parallel. It works by analyzing the
Fortran Intermediate Representation (FIR) produced by the LLVM Flang
compiler, running a five-phase pipeline that checks array subscript
patterns, detects scalar reductions, and applies a conservative
fallback. It produces a verdict (\textsc{Safe}, \textsc{Reduction}, or
\textsc{Unsafe}) along with an OpenMP directive hint and a full
explanation. On a 35-case test suite it achieves 100\% accuracy with
zero false positives --- it never tells you a loop is safe when it is
not.
\end{infobox}

\vfill
\begin{center}
{\color{navyblue}\rule{0.5\linewidth}{1pt}}\\[6pt]
{\small\color{darkgray}
Aryan Gupta \quad Arushi Vaidya \quad Arpita\\
Department of Computer Science and Engineering\\
RV College of Engineering, Bengaluru\\[4pt]
\url{https://github.com/AryanGupta21/FlangParallelAnalyzer}
}
\end{center}

\end{document}
